\section{Use past samples to define the current result}
A finite impulse response filter, or FIR filter, combines a bounded number of current
and past input samples. A coefficient-weighted contribution is called a tap. This
project uses eight taps. Unlike the preceding mixer's two simultaneous streams, the
filter reuses overlapping windows from one chronological stream.

The filter can change a signal's amplitude and phase according to its sample pattern.
Begin with an impulse and integer examples so that the output can be checked exactly.
The later listening fixture adds a 250 Hz triangular component to a specific 4 kHz
periodic component. An eight-sample average cancels the latter when each supplied
window contains complete periods. This arithmetic kernel neither opens a WAV file nor
resamples audio. A later project handles the file container.

The scalar and MSA paths implement the same exact arithmetic. Instruction count and
memory-read work can then be compared without changing rounding, accumulation width,
or sample ownership. A shorter sequence that overflows a partial sum implements a
different filter on some allowed inputs.
